\documentclass[11pt]{letter}
\usepackage[margin=1in]{geometry}
\usepackage{hyperref}
\usepackage{xcolor}

\signature{Tatsuki Onishi, MD\\
Department of Anesthesiology\\
University Hospital, Japan\\
\texttt{bougtoir@gmail.com}}

\address{Tatsuki Onishi\\
Department of Anesthesiology\\
University Hospital, Japan\\
\texttt{bougtoir@gmail.com}}

\date{\today}

\begin{document}

\begin{letter}{%
Editors\\
Journal of Machine Learning Research}

\opening{Dear Editors,}

I am writing to submit our manuscript entitled
\textbf{``Spectral Causality: Causal Direction Estimation via Magnetic
Laplacians and Hodge Decomposition''}
for consideration as a regular article in the
\emph{Journal of Machine Learning Research}.

\paragraph{Summary.}
This paper introduces \emph{spectral causality}, a framework that
estimates causal directions among observed variables by exploiting the
spectral structure of the magnetic Laplacian.  Unlike classical
Laplacian methods, the magnetic Laplacian's complex-valued eigenvectors
encode edge directionality as phase, enabling causal direction
estimation from spectral structure alone.  Our principal contributions
are:

\begin{enumerate}
\item A \textbf{Directional Predictability Index (DPI)} that resolves
  the circularity in earlier utility-based causal formulations, with
  partial identifiability guarantees under the additive noise model and
  explicit convergence rates for each component.

\item A \textbf{Hodge-theoretic decomposition} of causal edge flows
  into gradient (DAG-compatible), curl (feedback), and harmonic
  components, yielding a gradient energy ratio
  $r_{\mathrm{gradient}}$ that serves as a quantitative DAG-adequacy
  diagnostic---a capability absent from existing causal discovery
  methods.

\item A \textbf{scale invariance theorem} establishing that causal
  structure emergence depends on knowledge \emph{quality} (sign
  pattern) rather than quantity (magnitude), together with a
  phase-transition analysis characterizing a U-shaped knowledge
  quality curve.

\item Comprehensive \textbf{multi-dataset validation} on synthetic
  random DAGs ($n{=}5$--$20$, $N{=}200$--$1000$), the Sachs protein
  signaling network ($n{=}11$, 17 ground-truth edges), and the UCI
  Heart Disease dataset ($n{=}5$, $N{=}297$), benchmarked against
  DirectLiNGAM, the PC algorithm, GES, and NOTEARS.
\end{enumerate}

\paragraph{Relevance to JMLR.}
The manuscript presents a new principled algorithm (spectral causality)
grounded in spectral graph theory and Hodge theory, with sound
empirical validation on both synthetic and real-world benchmarks.  The
theoretical analysis (17 theorem-like statements with formal proofs)
advances understanding of the relationship between spectral structure
and causal ordering.  We believe the work is of broad interest to the
machine learning community, as causal discovery is a central topic
spanning methodology, theory, and applications.

\paragraph{Originality.}
This work has not been published previously in any journal or
conference proceedings.  A preprint may be made available on arXiv
concurrent with this submission.

\paragraph{Disclosure.}
No external funding was received for this work.  The author declares
no conflicts of interest.

\paragraph{Suggested action editors.}
Based on their expertise in causal discovery and spectral methods, we
suggest the following action editors who may be suitable to handle this
submission:

\begin{itemize}
\item \textbf{Bernhard Sch\"olkopf} --- causal inference, kernel methods
\item \textbf{Jonas Peters} --- causal discovery, additive noise models
\item \textbf{Peter Spirtes} --- constraint-based causal discovery
\item \textbf{Kun Zhang} --- causal discovery methodology
\end{itemize}

Thank you for considering this manuscript.  I look forward to receiving
your editorial decision.

\closing{Sincerely,}

\end{letter}
\end{document}
